% ===================================================================
% CHAPTER 9: CONCLUSION AND DISCUSSION
% ===================================================================
\chapter{Заклучок и дискусија}
\label{chap:conclusion}

\section{Одговори на истражувачките прашања}

Главното прашање на трудот беше колкави се физичките трошоци на заборавањето засновано на
SISA наспроти наивното повторно тренирање, кога ресурсно ограничен федериран јазол на работ
мора да отстрани идентификувано компромитирани податоци. Кампањата од десет спарени
извршувања дава следните одговори.

\textbf{RQ1 --- пресметковен и енергетски трошок.} SISA го намалува времето до обновување од
медијана 242,1 s на 29,5 s (\textbf{8,2×}) и нето-потрошената енергија од 0,093 Wh на 0,017 Wh
(\textbf{5,4×}). Двете со $p = 0{,}0020$ --- долната граница при $N = 10$ --- и Cliff's
$\delta = +1{,}00$, односно целосно раздвојување: ниту едно извршување со SISA не се
преклопува со ниту едно наивно извршување. Разликата не е статистички суптилна; таа е
разлика меѓу четири минути и половина минута.

\textbf{RQ2 --- термичко и влезно-излезно однесување.} Термичкото оптоварување се намалува од
медијана 233 s на 27,5 s со придушување ($\delta = +1{,}00$). Одговорот, меѓутоа, бара
корекција на почетната претпоставка: \textbf{врвната температура не е показателот што раздвојува} --- уред без активно ладење се заситува кон истата граница под секое одржливо оптоварување, и
опсезите на двата крака се преклопуваат. Показателот е \textbf{траењето} на придушувањето.
Претпоставената деградација на тактот, пак, воопшто не се потврди ($p = 0{,}25$,
$\delta = -0{,}22$) и се известува како нулти наод.

Трошокот навистина се поместува кон складирањето: SISA запишува 12,05 MB наспроти 2,75 MB
во мерниот прозорец. Насоката на H4 е потврдена, но \textbf{апсолутната големина ја обезвреднува
нејзината важност} --- дванаесет мегабајти не оптоваруваат ниту A1-картичка, а времето на
чекање за влез/излез останува под еден процент во двата крака. Хипотезата беше формулирана
како тестирана, а не претпоставена, и токму затоа занемарливиот резултат може да се пријави
како резултат, а не како неуспех.

\textbf{RQ3 --- задржување на корисноста.} Одговорот има два слоја, и само вториот е конечен.

При прагот на одлучување што реално се распоредува (0,5), обновениот модел со SISA
\textbf{колабира во мнозинската класа}: предвидува дека секој тек е напад, со 0 до 2 точно
детектирани бенигни примери од околу 3\,494, избалансирана точност 0,500 --- случаен избор --- и
MCC 0,017, наспроти 0,598 и 0,177 за наивниот пристап. Прочитано преку одѕивот моделот
изгледа совршен (1,000 наспроти 0,960); прочитано преку целосната матрица на конфузија, тој
одѕив е тривијален.

Но колапсот е \textbf{калибрациски, не информациски}. Со праг подесен низ целиот опсег, SISA
достигнува избалансирана точност \textbf{0,709 наспроти 0,671} за наивниот пристап --- подобра во 9
од 10 почетни вредности, и подобра истовремено по специфичност (0,694 наспроти 0,662) и по
MCC (0,171 наспроти 0,142). Калибрацијата ѝ носи на SISA +0,209, а на наивниот пристап само
+0,073.

Одговорот на RQ3 затоа гласи: \textbf{SISA ја задржува корисноста --- дури и малку ја надминува
онаа на наивното повторно тренирање --- но само ако прагот се калибрира.} Со подразбираниот
праг таа испорачува модел што е формално одличен и практично бескорисен. Вредноста 0,709 е
горна граница, добиена со праг избран на истото множество на кое се известува
(\hyperref[chap:limitations]{Дел~\ref*{chap:limitations}}, §8.13).

\textbf{RQ4 --- вкупен трошок на сопственост.} Режискиот трошок на SISA изнесува 250 контролни
точки, 43,67 MB и медијана од 5,56 s запишување низ десет рунди обука. Медијаната на
заштеденото време при едно обновување е 213,7 s. Билансот е \textbf{над триесеткратен во полза на
SISA}, дури и ако се плати целата цена на подготовката, а барањето за заборавање пристигне
само еднаш.

\section{Синтеза: што всушност е измерено}

Ако четирите одговори се спојат, сликата е поинаква и од почетната хипотеза, и од првиот
впечаток што го остава оценувањето при подразбираниот праг.

\begin{quote}
\textbf{SISA го прави заборавањето осумкратно побрзо, петкратно поенергетски ефикасно и
осумкратно помалку термички оптоварувачко --- без загуба во корисноста, под услов прагот на
одлучување да биде калибриран. Некалибрирана, таа испорачува модел што е формално одличен
и практично бескорисен.}
\end{quote}

Двата дела од таа реченица се измерени, не претпоставени: заштедата има целосно раздвојување
низ десет спарени извршувања, а условот за калибрација е квантифициран во одговорот на RQ3
погоре.

Клучно за толкувањето е што колапсот \textbf{не доаѓа од заборавањето}. Моделите од Фаза 1 веќе
имаат избалансирана точност 0,500, значи тој му претходи; а контролна проверка врз составните
модели го локализира местото --- поединечно тие препознаваат бенигни текови (избалансирана
точност 0,545–0,585), додека параметарската средна вредност на нивните тежини не препознава
ниту еден (\hyperref[chap:results]{Дел~\ref*{chap:results}}, §7.5). Дефектот затоа е својство на \textbf{начинот на вклопување} на
SISA во федерацијата --- на документираното отстапување во кое клиентот испраќа средна вредност
на тежините наместо ансамбл од предвидувања. Измерената цена на заборавањето е измерена
коректно.

Оттука и двојната поука: краткорочно, калибриран праг е задолжителен дел од распоредувањето;
долгорочно, агрегација што ги чува составните модели одвоени го отстранува дефектот во
корен.

\section{Дискусија: предвидени приговори}

Дизајнот на трудот е насочен кон неколку предвидливи приговори; тие се изнесуваат и се
одговараат овде експлицитно.

\textbf{„Метриката е бесмислена при нерамнотежа од 96,5\,\% --- одѕив 1,0 е гласање на мнозинството.“}
Приговорот е \textbf{точен}, и токму затоа секој зачуван глобален модел е повторно оценет врз
издвоено тест-множество со целосна матрица на конфузија. Заклучоците се засноваат на
избалансираната точност, специфичноста и MCC, не на одѕивот и F1. Тоа изнесе наод \textbf{против}
предметот на истражувањето --- колапс во мнозинската класа при подразбираниот праг --- и тој е
известен како таков, пред да се покаже дека е калибрациски. Овој приговор би бил фатален само
доколку одѕивот и F1 беа известени неквалификувано.

\textbf{„SISA троши помалку затоа што прави помалку --- споредбата е кружна.“} Првиот дел од
приговорот е \textbf{точен, и е измерен}: односот на извршената работа е 8,44×, односот на времето
8,23×, а пропусноста по ред е практично иста во двата крака (0,98×). Забрзувањето навистина е
избегнатата пресметка, ниту повеќе ниту помалку.

Тоа сепак не ја прави споредбата кружна, поради две причини. Прво, двата крака завршуваат со
\textbf{иста гаранција} --- егзактно локално заборавање на истото множество примери --- па се мери
цената на постигнување на фиксна цел, а не цената на различни количества работа. Второ, фактот
дека забрзувањето е точно еднакво на избегнатата пресметка е она што го прави \textbf{предвидливо}:
тоа не е случајна последица на имплементацијата, туку структурна последица на $S$, $R$ и
положбата на компромитираниот сегмент, па може да се пресмета однапред --- и, како што покажува
§7.4, може да се сведе на нула кога компромитирањето не е локално.

Треба, сепак, да се додаде и трета квалификација: целта е \emph{отстранување}, не \emph{квалитет на
моделот}. Добиениот модел со SISA е значително полошо калибриран, и при подразбираниот праг
тоа го прави бескорисен. Заштедата е бесплатна само откако прагот ќе се подеси --- што чини
ништо, но мора да се направи свесно.

\textbf{„Дали преобуката од нула е праведна референтна основа?“} Таа е избрана затоа што е
златниот стандард за егзактно заборавање --- единствената постапка што по дефиниција ја
задоволува равенката на егзактност. Поевтини алтернативи (дообучување, помал број епохи) би
дале поповолна споредба за наивниот крак, но повеќе не би гарантирале егзактност, па
споредбата би престанала да биде споредба на иста гаранција. Буџетот е изедначен: десет
епохи врз задржаните податоци, што е точно локалниот буџет од Фаза 1.

\textbf{„Предноста зависи од бројот на рунди.“} Точно, и се известува како наод за скалирањето,
не како недостаток. Враќањето наназад го носи труењето нанапред низ сите рунди, па SISA го
повторува засегнатиот шард за сите десет; предноста се намалува со повеќе рунди и расте со
помалку.

\textbf{„Ограниченото труење нема глобален ефект --- зошто тогаш обновување?“} Затоа што
отстранувањето е наложено од \textbf{откривање или од регулатива}, не од измерена штета. Тоа е
рамката поставена уште во \hyperref[chap:problem]{Дел~\ref*{chap:problem}}: прашањето не е колку штети нападот, туку колку
чини отстранувањето. Отпорноста на федерацијата кон просторно ограничено труење е самостоен
наод и е известена како таков.

\textbf{„Еден уред --- колку се генерализира?“} Не многу, и тоа е наведено во \hyperref[chap:limitations]{Дел~\ref*{chap:limitations}}.
Трудот е студија на случај врз еден уред со целосно обелоденет хардвер. Придонесот не е
универзална константа, туку \textbf{методологија на мерење} и ред на големина за еден реален,
претставителен уред.

\section{Инженерски лекции}

Спроведувањето на физички експеримент откри низа замки што не се видливи во симулација.
Тие се наведуваат затоа што претставуваат самостоен придонес за секого што сака да мери
работни оптоварувања на уреди на работ.

\begin{itemize}
  \item \textbf{Испораката на струја е скриен збунувачки фактор.} Тенок кабел меѓу мерачот и уредот паѓа доволно напон за да го подцени напојувањето и да предизвика придушување што нема никаква врска со работното оптоварување --- и тоа додека самиот мерач покажува уредни 5,1 V. Проблемот беше откриен во пилотот и ублажен со кабел од 5 A и фиксирање на преговарањето за напојување во постојаната меморија на уредот. Како што покажува \hyperref[chap:limitations]{Дел~\ref*{chap:limitations}}, §8.4, тој не е целосно отстранет --- што е причина повеќе секое извршување да го проверува регистарот за придушување пред почетокот.
  \item \textbf{Прагот за ладење мора да биде релативен, не апсолутен.} Првичниот услов да се чека температура под 40 °C е физички недостижен на уред без вентилатор, чиј праг во мирување е 60–79 °C во зависност од просторијата. Замената со услов за \textbf{стабилно плато} --- опсег до 2 °C во текот на 120 s --- ја постигнува вистинската цел: секое извршување да започне од иста состојба, а не од иста бројка.
  \item \textbf{Лепливите битови за придушување се замка.} Битовите што бележат дека настанот се случил некогаш остануваат поставени до рестарт. Броење на сè што е различно од нула драстично го преувеличува придушувањето; се бројат само живите битови.
  \item \textbf{Телеметријата преку SSH мора да се подигне како услуга.} Процес стартуван со \texttt{nohup} или \texttt{setsid} преку SSH не опстојува поуздано; потребно е привремена системска единица.
  \item \textbf{Виртуелната меморија се враќа сама.} На Debian 13 новиот механизам ја обновува по рестарт, тивко, и ја загадува токму метриката за влезно-излезни операции што се мери. Оневозможувањето мора да биде трајно запишано.
  \item \textbf{Почетната вредност на кешот мора да одговара на почетната вредност при извршувањето.} Труењето е вградено во кешот на јазолот според распоредувањето по шардови изведено од почетната вредност; истото распоредување се реконструира при извршувањето. Несовпаѓање ги распрснува фиксираните редови во погрешни сегменти и ја обесмислува постапката. Проверката е едноставна и треба да биде задолжителна: секое исправно подготвено извршување мора да пријави \texttt{\{шард 1: сегмент 3\}}.
\end{itemize}

\section{Придонеси}

Трудот дава четири придонеси:

\begin{enumerate}
  \item \textbf{Емпириско мерење на цената на заборавањето врз вистински хардвер на работ}, изразено во секунди, ват-часови, секунди со термичко придушување и запишани мегабајти --- величини што постојната литература, која мери во епохи и рунди во симулација, не ги дава.
  \item \textbf{Претходно регистриран спарен експериментален дизајн} со непараметриска статистика, големини на ефектот и bootstrap интервали, спроведен низ десет спарени извршувања со балансиран редослед, со објавени податоци по извршување.
  \item \textbf{Наод што го коригира сопственото очекување}: демонстрација дека при изразена класна нерамнотежа стандардните показатели за системи за откривање упади ја \textbf{превртуваат} вистината. Оценета со одѕив и F1, SISA изгледа совршена; оценета со избалансирана точност при подразбираниот праг, изгледа бескорисна; оценета со калибриран праг, е подобра од наивниот пристап. Сите три читања се на \textbf{истиот модел}. Наодот е добиен против предметот на истражувањето и покажува дека изборот на показател и на праг не е техничка ситница, туку го определува заклучокот.
  \item \textbf{Методологија и инженерски лекции} за мерење физички трошоци на уреди на работ, од изолацијата на влезно-излезниот шум до изборот на показател за термичко оптоварување.
\end{enumerate}

Конечниот заклучок е условен, а не победнички, и тоа е негова предност: \textbf{SISA е убедливо
поевтиниот начин да се заборави на уред на работ, а нејзиниот трошок на подготовка е
занемарлив; корисноста не е жртвувана, но се наплатува ако прагот на одлучување се остави
некалибриран. Поправката е бесплатна, а трајното решение --- агрегација што ги чува составните
модели одвоени --- останува отворена задача.}

